% ============================================================
%  Bolometer Simulator --- Developer Guide
%  Compile: pdflatex bolometer_devguide.tex  (twice for TOC)
% ============================================================
\documentclass[11pt,a4paper]{article}

% ---- packages ------------------------------------------------
\usepackage[T1]{fontenc}
\usepackage[utf8]{inputenc}
\usepackage[margin=2.5cm]{geometry}
\usepackage{amsmath,amssymb}
\usepackage{booktabs}
\usepackage{xcolor}
\usepackage{listings}
\usepackage{hyperref}
\usepackage{graphicx}
\usepackage{float}
\usepackage{enumitem}
\usepackage{titlesec}
\usepackage{fancyhdr}
\usepackage{microtype}
\usepackage{mdframed}
\usepackage{parskip}
\usepackage{textcomp}

% ---- colours -------------------------------------------------
\definecolor{matlabblue}{HTML}{0072BD}
\definecolor{matlabgreen}{HTML}{217821}
\definecolor{matlabgray}{HTML}{808080}
\definecolor{matlabcomment}{HTML}{3D7A14}
\definecolor{matlabstring}{HTML}{A2101F}
\definecolor{codebg}{HTML}{F5F7FA}
\definecolor{rulecolor}{HTML}{CDD3DC}
\definecolor{titleblue}{HTML}{0D4B8C}
\definecolor{noteblue}{HTML}{EBF3FB}
\definecolor{noteborder}{HTML}{2E86C1}
\definecolor{warnbg}{HTML}{FEF9E7}
\definecolor{warnborder}{HTML}{D4AC0D}

% ---- MATLAB listing style -----------------------------------
\lstdefinestyle{matlab}{
  language=Matlab,
  backgroundcolor=\color{codebg},
  basicstyle=\ttfamily\footnotesize,
  keywordstyle=\color{matlabblue}\bfseries,
  commentstyle=\color{matlabcomment}\itshape,
  stringstyle=\color{matlabstring},
  numberstyle=\tiny\color{matlabgray},
  numbers=left,
  numbersep=8pt,
  stepnumber=1,
  frame=single,
  rulecolor=\color{rulecolor},
  framesep=4pt,
  breaklines=true,
  breakatwhitespace=false,
  showstringspaces=false,
  tabsize=4,
  captionpos=b,
  xleftmargin=12pt,
  xrightmargin=4pt,
  morekeywords={function,end,if,for,while,true,false,nargin,fieldnames,isfield,
                logspace,linspace,randn,sqrt,abs,mean,std,log10,max,min,
                sign,zeros,ones,repmat,strcmp,any,all,sum,tinv,reshape,size,
                length,fprintf,subplot,hold,grid,box,plot,semilogx,semilogy,
                loglog,errorbar,fill,yline,xline,xlim,ylim,set,title,xlabel,
                ylabel,legend,sgtitle,figure,exportgraphics,linkaxes,flipud,
                fieldnames,mat2str,sprintf},
}
\lstset{style=matlab}

% ---- callout environments -----------------------------------
\newmdenv[
  backgroundcolor=noteblue,
  linecolor=noteborder,
  linewidth=2pt,
  topline=false, rightline=false, bottomline=false,
  innerleftmargin=10pt, innerrightmargin=8pt,
  innertopmargin=6pt, innerbottommargin=6pt,
]{infobox}

\newmdenv[
  backgroundcolor=warnbg,
  linecolor=warnborder,
  linewidth=2pt,
  topline=false, rightline=false, bottomline=false,
  innerleftmargin=10pt, innerrightmargin=8pt,
  innertopmargin=6pt, innerbottommargin=6pt,
]{warnbox}

% ---- headers ------------------------------------------------
\pagestyle{fancy}
\fancyhf{}
\lhead{\textcolor{matlabgray}{\small Bolometer Simulator --- Developer Guide}}
\rhead{\textcolor{matlabgray}{\small\thepage}}
\renewcommand{\headrulewidth}{0.4pt}

% ---- section formatting -------------------------------------
\titleformat{\section}{\large\bfseries\color{titleblue}}{\thesection}{0.8em}{}
\titleformat{\subsection}{\normalsize\bfseries\color{titleblue}}{\thesubsection}{0.7em}{}
\titleformat{\subsubsection}{\normalsize\bfseries}{\thesubsubsection}{0.6em}{}

% ---- hyperref -----------------------------------------------
\hypersetup{
  colorlinks=true,
  linkcolor=titleblue,
  urlcolor=matlabblue,
  citecolor=matlabblue,
  pdftitle={Bolometer Simulator --- Developer Guide},
  pdfauthor={gelfprog},
}

% ============================================================
\begin{document}

% ---- title page ---------------------------------------------
\begin{titlepage}
  \centering
  \vspace*{2cm}
  {\Huge\bfseries\color{titleblue} Bolometer Simulator\\[0.4em]}
  {\LARGE Developer Guide\par}
  \vspace{1.5cm}
  {\large\ttfamily bolometer/ --- MATLAB codebase\par}
  \vspace{0.8cm}
  \rule{12cm}{0.4pt}\\[0.6cm]
  {\large
    \textbf{Repository:} \texttt{github.com/AltFed/gelfprog}\\[0.3em]
    \textbf{Corso:} Diagnostiche per plasmi\\[0.3em]
    \textbf{Docenti:} Dr.~A.~Murari \& Prof.~M.~Gelfusa\\[0.3em]
    \textbf{Data:} Maggio 2026
  }
  \vspace{1.5cm}

  \begin{mdframed}[backgroundcolor=codebg,linecolor=rulecolor]
  \centering\ttfamily\large
  \textcolor{matlabblue}{P} $\to$
  \textcolor{matlabgreen}{$\Delta$T} $\to$
  \textcolor{matlabgreen}{$\Delta$R} $\to$
  \textcolor{matlabgreen}{$\Delta$V} $\to$
  \textcolor{matlabstring}{+\,noise} $\to$
  \textcolor{matlabblue}{P$_\text{rec}$}
  \end{mdframed}
  \vfill
  {\small\color{matlabgray}
    Questo documento è destinato agli sviluppatori che vogliono comprendere,
    modificare o estendere il simulatore. Contiene l'architettura completa,
    le equazioni fisiche, i listati commentati e le istruzioni di estensione.}
\end{titlepage}

\tableofcontents
\newpage

% ============================================================
\section{Panoramica del Progetto}

Il simulatore implementa la catena di misura completa di un bolometro resistivo
a foglio metallico (tipo gold-foil usato su JET, FTU, ITER). La catena è:

\begin{equation}
  P \;\xrightarrow{\text{bilancio termico}}\;
  \Delta T \;\xrightarrow{\text{TCR}}\;
  \Delta R \;\xrightarrow{\text{polarizzazione}}\;
  \Delta V_{\text{ideal}}
  \;\xrightarrow{\text{+rumore}}\;
  \Delta V_{\text{noisy}}
  \;\xrightarrow{\text{inversione}}\;
  P_{\text{rec}}
\end{equation}

\subsection*{Struttura del repository}

\begin{verbatim}
gelfprog/
+-- bolometer/
|   +-- bolometer_params.m        parametri fisici e di rumore
|   +-- forward_model.m           P -> DeltaT -> DeltaR -> DeltaV  (forward model)
|   +-- add_noise_to_signal.m     DeltaV + rumore gaussiano (N realizzazioni MC)
|   +-- reconstruct_power.m       DeltaV_noisy -> P_rec  (inversione + statistiche)
|   +-- run_power_sweep.m         sweep logaritmico + figura 3 pannelli
|   +-- sensitivity_analysis.m    sweep di R0,alpha,G,I_bias -> S/NEP/P_sat/DR
|   +-- main.m                    script master configurabile
|   +-- CONTEXT.md                contesto aggiornato ad ogni modifica
|   +-- test_block1..5.m          smoke-test per ogni blocco
|   +-- *.png                     figure generate
+-- bolometer_simulator.pptx      presentazione (10 slide)
+-- bolometer_devguide.pdf        questo documento
\end{verbatim}

\subsection*{Come eseguire il simulatore}

Aprire MATLAB, spostarsi nella cartella \texttt{bolometer/} ed eseguire:

\begin{lstlisting}[caption={Esecuzione rapida con parametri default}]
cd bolometer/
main
\end{lstlisting}

Per personalizzare i parametri fisici o di rumore, editare la sezione
\texttt{CUSTOM\_PARAMS} in \texttt{main.m}:

\begin{lstlisting}[caption={Override parametri in main.m}]
CUSTOM_PARAMS = {'R0', 200, 'G', 5e-5, 'I_bias', 2e-3};
\end{lstlisting}

% ============================================================
\section{Fondamenti Fisici}

\subsection{Modello termico in regime stazionario}

L'equazione di bilancio termico del bolometro è:

\begin{equation}
  C\,\frac{d(\Delta T)}{dt} + G\,\Delta T = P(t)
\end{equation}

dove $C$~[J/K] è la capacità termica e $G$~[W/K] la conduttanza termica.
In regime stazionario ($d/dt = 0$):

\begin{equation}
  \boxed{\Delta T = \frac{P}{G}}
  \label{eq:deltaT}
\end{equation}

La costante di tempo termica è $\tau = C/G$. Per i parametri default
$\tau = 10\,\mathrm{ms}$.

\subsection{Effetto TCR (Temperature Coefficient of Resistance)}

La variazione di resistenza del sensore è proporzionale alla variazione
di temperatura tramite il coefficiente TCR $\alpha$~[1/K]:

\begin{equation}
  \Delta R = R_0\,\alpha\,\Delta T
  \label{eq:deltaR}
\end{equation}

Per l'oro $\alpha \approx 3.4 \times 10^{-3}$~1/K, per il platino
$\alpha \approx 3.9 \times 10^{-3}$~1/K.

\subsection{Lettura con corrente di polarizzazione}

Con una corrente di polarizzazione costante $I_{\text{bias}}$:

\begin{equation}
  \Delta V = I_{\text{bias}}\,\Delta R
  \label{eq:deltaV}
\end{equation}

Combinando le equazioni \eqref{eq:deltaT}--\eqref{eq:deltaV}:

\begin{equation}
  \boxed{\Delta V = \underbrace{\frac{I_{\text{bias}}\,R_0\,\alpha}{G}}_{S}\cdot P
  \qquad \Rightarrow \qquad S = \frac{\Delta V}{P} \quad [\text{V/W}]}
  \label{eq:sensitivity}
\end{equation}

\subsection{Saturazione}

L'uscita satura in due modi:
\begin{itemize}[noitemsep]
  \item \textbf{Termica:} $\Delta T > \Delta T_{\max}$ (danneggiamento sensore)
  \item \textbf{Elettronica:} $|\Delta V| > V_{\text{sat}}$ (amplificatore al rail)
\end{itemize}
La potenza di saturazione è:
\begin{equation}
  P_{\text{sat}} = \frac{V_{\text{sat}}}{S}
\end{equation}

\subsection{Modello del rumore}

Tre sorgenti indipendenti, sommate in quadratura:

\begin{align}
  V_J &= \sqrt{4\,k_B\,T_0\,R_0\,\Delta f}  && \text{(Johnson/termico)} \\
  V_A &= e_n\,\sqrt{\Delta f}               && \text{(amplificatore)} \\
  V_Q &= \frac{V_{\text{range}}}{2^N\,\sqrt{12}} && \text{(quantizzazione ADC)} \\[6pt]
  V_{\text{noise}} &= \sqrt{V_J^2 + V_A^2 + V_Q^2}
\end{align}

Il \textbf{Noise Equivalent Power} è:
\begin{equation}
  \text{NEP} = \frac{V_{\text{noise}}}{S}
\end{equation}

\subsubsection*{Dynamic Range}

\begin{equation}
  \text{DR} = \frac{P_{\text{sat}}}{\text{NEP}}
  = \frac{V_{\text{sat}}/S}{V_{\text{noise}}/S}
  = \frac{V_{\text{sat}}}{V_{\text{noise}}}
  \label{eq:DR}
\end{equation}

\begin{infobox}
\textbf{Risultato chiave:} il DR è \emph{indipendente} da $S$.
Variare $R_0$, $\alpha$, $G$ o $I_{\text{bias}}$ sposta in modo uguale
$P_{\text{sat}}$ e NEP, lasciando il DR invariato.
Con i parametri default: $\text{DR} \approx 101$~dB.
Per aumentare il DR occorre agire sull'ADC (più bit o range ridotto)
o sull'amplificatore ($e_n$ minore).
\end{infobox}

% ============================================================
\section{Documentazione dei Moduli}

Ogni modulo è un file MATLAB autonomo che riceve un parametri-struct
prodotto da \texttt{bolometer\_params} e restituisce uno struct di risultati.
Nessun modulo usa variabili globali.

% ------------------------------------------------------------
\subsection{\texttt{bolometer\_params.m} --- Parametri}

\subsubsection*{Scopo}
Costruisce e restituisce uno struct con tutti i parametri del simulatore.
Supporta l'override selettivo tramite coppie nome-valore.

\subsubsection*{Firma}
\begin{lstlisting}
params = bolometer_params()
params = bolometer_params('R0', 200, 'G', 5e-5)
\end{lstlisting}

\subsubsection*{Parametri dello struct}

\begin{center}
\begin{tabular}{llrl}
\toprule
Campo & Simbolo & Default & Unità \\
\midrule
\texttt{R0}        & $R_0$             & 100        & $\Omega$ \\
\texttt{alpha}     & $\alpha$          & $3.9\times10^{-3}$ & 1/K \\
\texttt{G}         & $G$               & $10^{-4}$  & W/K \\
\texttt{C}         & $C$               & $10^{-6}$  & J/K \\
\texttt{I\_bias}   & $I_{\text{bias}}$ & $10^{-3}$  & A \\
\texttt{T0}        & $T_0$             & 300        & K \\
\texttt{V\_sat}    & $V_{\text{sat}}$  & 1.0        & V \\
\texttt{dT\_max}   & $\Delta T_{\max}$ & 50         & K \\
\texttt{BW}        & $\Delta f$        & $10^3$     & Hz \\
\texttt{e\_n}      & $e_n$             & $10^{-8}$  & V/$\sqrt{\text{Hz}}$ \\
\texttt{N\_bits}   & $N$               & 16         & bit \\
\texttt{V\_ADC\_range} & $V_{\text{range}}$ & 2.0  & V \\
\texttt{N\_MC}     & ---                 & 1000       & --- \\
\midrule
\texttt{tau} & $\tau=C/G$ & $10^{-2}$ & s \\
\texttt{S}   & $S = I_{\text{bias}}R_0\alpha/G$ & 3.9 & V/W \\
\bottomrule
\end{tabular}
\end{center}

\subsubsection*{Implementazione chiave}

\begin{lstlisting}[caption={Override e ricalcolo campi derivati in bolometer\_params.m}]
% override with name-value pairs
for k = 1:2:length(varargin)
    field = varargin{k};
    value = varargin{k+1};
    if ~isfield(params, field)
        error('bolometer_params: unknown field "%s"', field);
    end
    params.(field) = value;
end

% recompute derived quantities after any override
params.tau = params.C / params.G;
params.S   = params.I_bias * params.R0 * params.alpha / params.G;
\end{lstlisting}

\begin{warnbox}
\textbf{Attenzione:} i campi \texttt{tau} e \texttt{S} sono ricalcolati
automaticamente dopo ogni override. Non modificarli a mano nello struct.
\end{warnbox}

% ------------------------------------------------------------
\subsection{\texttt{forward\_model.m} --- Modello Diretto}

\subsubsection*{Scopo}
Implementa la catena fisica $P \to \Delta T \to \Delta R \to \Delta V$
inclusa la saturazione.

\subsubsection*{Firma}
\begin{lstlisting}
result = forward_model(P, params)
\end{lstlisting}

\subsubsection*{Input / Output}

\begin{center}
\begin{tabular}{lll}
\toprule
Campo output & Tipo & Descrizione \\
\midrule
\texttt{DeltaT}       & array & incremento di temperatura $[K]$ \\
\texttt{DeltaR}       & array & variazione di resistenza $[\Omega]$ \\
\texttt{DeltaV\_ideal}& array & tensione ideale (senza saturazione) $[V]$ \\
\texttt{DeltaV}       & array & tensione clippata a $V_{\text{sat}}$ $[V]$ \\
\texttt{is\_saturated}& logical & maschera punti saturi \\
\texttt{P\_sat}       & scalare & potenza di saturazione $[W]$ \\
\texttt{S}            & scalare & sensibilità usata $[V/W]$ \\
\bottomrule
\end{tabular}
\end{center}

\subsubsection*{Implementazione}

\begin{lstlisting}[caption={forward\_model.m --- corpo principale}]
% step 1: thermal balance (steady-state)
result.DeltaT = P / params.G;

% step 2: resistance change via TCR
result.DeltaR = params.R0 * params.alpha .* result.DeltaT;

% step 3: ideal voltage (current-bias readout)
result.DeltaV_ideal = params.I_bias .* result.DeltaR;

% step 4: saturation --- two mechanisms
thermal_sat  = result.DeltaT  > params.dT_max;
electric_sat = abs(result.DeltaV_ideal) > params.V_sat;
result.is_saturated = thermal_sat | electric_sat;

result.DeltaV = result.DeltaV_ideal;
result.DeltaV(result.is_saturated) = ...
    sign(result.DeltaV_ideal(result.is_saturated)) * params.V_sat;

result.S     = params.S;
result.P_sat = params.V_sat / params.S;
\end{lstlisting}

% ------------------------------------------------------------
\subsection{\texttt{add\_noise\_to\_signal.m} --- Modulo Rumore}

\subsubsection*{Scopo}
Aggiunge rumore realistico al segnale ideale $\Delta V$.
Supporta la generazione di $N$ realizzazioni Monte Carlo in parallelo
tramite operazioni vettorizzate su matrice $[M \times N]$.

\subsubsection*{Firma}
\begin{lstlisting}
result = add_noise_to_signal(DeltaV, params)
result = add_noise_to_signal(DeltaV, params, N_realizations)
\end{lstlisting}

\subsubsection*{Output principali}

\begin{center}
\begin{tabular}{lll}
\toprule
Campo & Dimensioni & Descrizione \\
\midrule
\texttt{DeltaV\_noisy}  & $[M \times N]$ & segnale rumoroso \\
\texttt{V\_J}           & scalare & Johnson noise rms $[V]$ \\
\texttt{V\_A}           & scalare & amplifier noise rms $[V]$ \\
\texttt{V\_Q}           & scalare & ADC quant.\ noise rms $[V]$ \\
\texttt{V\_noise\_total}& scalare & rumore totale $[V]$ \\
\texttt{SNR}            & $[M \times 1]$ & SNR per punto \\
\texttt{NEP}            & scalare & noise equivalent power $[W]$ \\
\bottomrule
\end{tabular}
\end{center}

\subsubsection*{Implementazione}

\begin{lstlisting}[caption={add\_noise\_to\_signal.m --- calcolo rumore e generazione MC}]
kB = 1.380649e-23;

% 1. Johnson noise (thermal agitation in the resistor)
V_J = sqrt(4 * kB * params.T0 * params.R0 * params.BW);

% 2. Amplifier noise
V_A = params.e_n * sqrt(params.BW);

% 3. ADC quantization noise
LSB = params.V_ADC_range / 2^params.N_bits;
V_Q = LSB / sqrt(12);

% Total noise (independent sources -> quadrature sum)
V_noise = sqrt(V_J^2 + V_A^2 + V_Q^2);

% Generate M x N_realizations Gaussian noise matrix
DeltaV = DeltaV(:);          % ensure column vector
M = length(DeltaV);
noise_matrix = V_noise * randn(M, N_realizations);
DeltaV_noisy = DeltaV + noise_matrix;

% Hard clip at ADC input range
V_half = params.V_ADC_range / 2;
DeltaV_noisy = max(-V_half, min(V_half, DeltaV_noisy));
\end{lstlisting}

\begin{infobox}
\textbf{Nota sulla vettorizzazione:} la generazione MC è completamente
vettorizzata. Per $M = 120$ punti e $N = 1000$ realizzazioni, la matrice
\texttt{noise\_matrix} è $120 \times 1000$ e viene creata in un'unica
chiamata a \texttt{randn}. Nessun ciclo \texttt{for}.
\end{infobox}

% ------------------------------------------------------------
\subsection{\texttt{reconstruct\_power.m} --- Ricostruzione Inversa}

\subsubsection*{Scopo}
Inverte la catena: $\Delta V_{\text{noisy}} \to P_{\text{rec}}$.
Calcola statistiche Monte Carlo e classifica ogni punto in uno dei
tre regimi (\texttt{noise}/\texttt{ok}/\texttt{saturated}).

\subsubsection*{Firma}
\begin{lstlisting}
result = reconstruct_power(DeltaV_noisy, params)
result = reconstruct_power(DeltaV_noisy, params, P_true)
\end{lstlisting}

Se \texttt{P\_true} è fornito, vengono calcolati anche RMSE ed errore relativo.

\subsubsection*{Output principali}

\begin{center}
\begin{tabular}{lll}
\toprule
Campo & Dimensioni & Descrizione \\
\midrule
\texttt{P\_mean}       & $[M \times 1]$ & media MC di $P_{\text{rec}}$ $[W]$ \\
\texttt{P\_std}        & $[M \times 1]$ & deviazione standard MC $[W]$ \\
\texttt{P\_ci95}       & $[M \times 1]$ & semi-ampiezza CI 95\% $[W]$ \\
\texttt{P\_all}        & $[M \times N]$ & tutte le realizzazioni $[W]$ \\
\texttt{is\_saturated} & $[M \times 1]$ & maschera punti saturi \\
\texttt{is\_noise\_dom}& $[M \times 1]$ & maschera punti noise-dominated \\
\texttt{regime}        & cell $[M\times1]$ & \texttt{'noise'/'ok'/'saturated'} \\
\texttt{NEP}           & scalare & NEP calcolato dai params $[W]$ \\
\texttt{RMSE}          & scalare & (se P\_true fornito) $[W]$ \\
\bottomrule
\end{tabular}
\end{center}

\subsubsection*{Implementazione: inversione e statistiche}

\begin{lstlisting}[caption={reconstruct\_power.m --- inversione e classificazione}]
% Core inversion: P_rec = DeltaV / S
P_all  = DeltaV_noisy / S;       % [M x N]
P_mean = mean(P_all, 2);          % [M x 1]
P_std  = std(P_all,  0, 2);       % [M x 1]

% 95% CI using t-distribution for finite N
t95    = tinv(0.975, N - 1);
P_ci95 = t95 * P_std / sqrt(N);

% Regime classification
sat_thresh      = 0.99 * params.V_sat;
mean_abs_DeltaV = abs(mean(DeltaV_noisy, 2));

is_saturated = mean_abs_DeltaV >= sat_thresh;
is_noise_dom = mean_abs_DeltaV <  V_noise;

regime = repmat({'ok'}, M, 1);
regime(is_noise_dom)  = {'noise'};
regime(is_saturated)  = {'saturated'};
\end{lstlisting}

\subsubsection*{Verifica analitica}

La deviazione standard delle realizzazioni MC deve approssimare il NEP:
\begin{equation}
  \sigma_{P_{\text{rec}}} = \frac{V_{\text{noise}}}{S} = \text{NEP}
\end{equation}
Questo è verificato numericamente nel test \texttt{test\_block3.m}:
$\sigma_P = 2.298 \times 10^{-6}$~W rispetto a NEP $= 2.260 \times 10^{-6}$~W.

% ------------------------------------------------------------
\subsection{\texttt{run\_power\_sweep.m} --- Power Sweep}

\subsubsection*{Scopo}
Esegue la pipeline completa su un range logaritmico di potenza
e produce una figura a 3 pannelli.

\subsubsection*{Firma}
\begin{lstlisting}
sweep = run_power_sweep()
sweep = run_power_sweep(params, options)
\end{lstlisting}

\subsubsection*{Opzioni (\texttt{options} struct)}

\begin{center}
\begin{tabular}{lrl}
\toprule
Campo & Default & Descrizione \\
\midrule
\texttt{P\_min}    & $10^{-9}$~W & potenza minima \\
\texttt{P\_max}    & $10^{0}$~W  & potenza massima \\
\texttt{N\_points} & 100         & punti nel range logaritmico \\
\texttt{N\_MC}     & 1000        & realizzazioni Monte Carlo \\
\texttt{do\_plot}  & \texttt{true} & genera la figura \\
\texttt{save\_fig} & \texttt{false} & salva la figura su disco \\
\texttt{fig\_path} & \texttt{'bolometer\_sweep.png'} & percorso output \\
\bottomrule
\end{tabular}
\end{center}

\subsubsection*{Struttura della figura (3 pannelli)}

\begin{enumerate}[noitemsep]
  \item \textbf{Pannello 1 (log-log):} $P_{\text{rec}}$ vs $P_{\text{true}}$,
        con codice colore per regime, banda CI95, linee NEP e $P_{\text{sat}}$
  \item \textbf{Pannello 2 (log-log):} SNR vs $P_{\text{true}}$, linea SNR=1
  \item \textbf{Pannello 3 (log-lin):} Curva di trasferimento $\Delta V(P)$:
        ideale vs clippata, $V_{\text{noise}}$ e $V_{\text{sat}}$
\end{enumerate}

\subsubsection*{Pattern di output sweep}

\begin{lstlisting}[caption={Esempio di output testuale (120 punti, N\_MC=1000)}]
NEP   = 2.2604e-06 W
P_sat = 2.5641e-01 W

Noise points  = 45 / 120    (P < NEP)
OK points     = 67 / 120
Sat. points   =  8 / 120    (P > P_sat)
RMSE (ok)     = 7.42e-08 W
\end{lstlisting}

% ------------------------------------------------------------
\subsection{\texttt{sensitivity\_analysis.m} --- Analisi di Sensibilità}

\subsubsection*{Scopo}
Varia un parametro alla volta su un range logaritmico
($\times 0.05$ ÷ $\times 20$ dal valore base) e calcola le metriche
$S$, NEP, $P_{\text{sat}}$, DR per ogni valore.

\subsubsection*{Firma}
\begin{lstlisting}
results = sensitivity_analysis()
results = sensitivity_analysis(params, options)
\end{lstlisting}

\subsubsection*{Logica del sweep}

\begin{lstlisting}[caption={Cuore di sensitivity\_analysis.m}]
for i = 1:N_params
    pname = sweep_defs(i).name;       % e.g. 'R0'
    for j = 1:N
        p = base_params;
        p.(pname) = sweep_defs(i).x_vals(j);
        % recompute derived fields
        p.tau = p.C / p.G;
        p.S   = p.I_bias * p.R0 * p.alpha / p.G;
        % compute noise and metrics
        V_n  = sqrt(V_J^2 + V_A^2 + V_Q^2);   % V_J depends on R0
        NEP  = V_n / p.S;
        Psat = p.V_sat / p.S;
        DR   = 20 * log10(Psat / NEP);         % = 20*log10(V_sat/V_n)
        ...
    end
end
\end{lstlisting}

\subsubsection*{Risultato chiave}

Come dimostrato in Eq.~\eqref{eq:DR}, DR = $V_{\text{sat}}/V_{\text{noise}}$:
\begin{itemize}[noitemsep]
  \item Per $\alpha$, $G$, $I_{\text{bias}}$: DR è \emph{esattamente} costante
        (le curve sono piatte al $\mu$dB)
  \item Per $R_0$: leggero calo a valori alti perché $V_J \propto \sqrt{R_0}$
        inizia a competere con $V_Q$, aumentando $V_{\text{noise}}$
  \item Variazione totale su range $\times 400$: $< 0.003$~dB
\end{itemize}

% ------------------------------------------------------------
\subsection{\texttt{main.m} --- Script Master}

\subsubsection*{Scopo}
Punto di ingresso unico. Definisce tutti i parametri nella sezione
\emph{USER SETTINGS} e chiama i moduli nelle sezioni abilitate dai flag
\texttt{RUN\_SWEEP} e \texttt{RUN\_SENS}.

\subsubsection*{Struttura interna}
\begin{lstlisting}[caption={Struttura di main.m}]
% 1. USER SETTINGS
CUSTOM_PARAMS = {};          % override parametri fisici
SWEEP.P_min   = 1e-9;        % sweep settings
SWEEP.N_MC    = 1000;
RUN_SWEEP = true;
RUN_SENS  = true;

% 2. INIT
p = bolometer_params(CUSTOM_PARAMS{:});
print_params(p);             % funzione locale in coda al file

% 3. POWER SWEEP
if RUN_SWEEP
    sweep = run_power_sweep(p, SWEEP);
    % ... stampa risultati
end

% 4. SENSITIVITY ANALYSIS
if RUN_SENS
    sens = sensitivity_analysis(p, SENS);
    % ... stampa noise budget
end

% 5. FINAL SUMMARY
% ... tabella con S, NEP, P_sat, DR
\end{lstlisting}

% ============================================================
\section{Flusso dei Dati tra Moduli}

\begin{verbatim}
bolometer_params()
      |
      v  params struct
  +---+-------------------------------------------+
  |                                               |
  v                                               v
forward_model(P, params)          sensitivity_analysis(params)
      |                                           |
      | .DeltaV_ideal                             | sweep S/NEP/DR
      v
add_noise_to_signal(DeltaV_ideal, params, N_MC)
      |
      | .DeltaV_noisy  [M x N_MC]
      v
reconstruct_power(DeltaV_noisy, params, P_true)
      |
      | .P_mean, .P_std, .P_ci95, .regime
      v
   [risultati + figure]
\end{verbatim}

Tutti i moduli condividono \emph{solo} il \texttt{params} struct.
Nessuna variabile globale. I risultati parziali fluiscono come
argomenti espliciti.

\begin{infobox}
\textbf{Interfaccia stabile:} \texttt{forward\_model}, \texttt{add\_noise\_to\_signal}
e \texttt{reconstruct\_power} possono essere chiamati indipendentemente.
La funzione \texttt{run\_power\_sweep} è una shell che li orchestra.
\end{infobox}

% ============================================================
\section{Risultati di Validazione Numerica}

\subsection{Consistenza del modello diretto}

\begin{center}
\begin{tabular}{rrrrr}
\toprule
$P$ [W] & $\Delta T$ [K] & $\Delta R$ [$\Omega$] & $\Delta V$ [V] & Saturato \\
\midrule
$10^{-9}$ & $10^{-5}$ & $3.9\times10^{-6}$ & $3.9\times10^{-9}$ & No \\
$10^{-6}$ & $10^{-2}$ & $3.9\times10^{-3}$ & $3.9\times10^{-6}$ & No \\
$10^{-3}$ & $10.0$    & $3.9$               & $3.9\times10^{-3}$ & No \\
$10^{-1}$ & $1000$    & ---                   & 1.000              & \textbf{Sì} \\
\bottomrule
\end{tabular}
\end{center}

\subsection{Budget del rumore (parametri default)}

\begin{center}
\begin{tabular}{lrr}
\toprule
Sorgente & $V_{\text{rms}}$ [V] & \% del totale \\
\midrule
Johnson $V_J$ & $4.07\times10^{-8}$ & 0.0\% \\
Amplificatore $V_A$ & $3.16\times10^{-7}$ & 0.1\% \\
ADC $V_Q$ (16-bit, 2V) & $8.81\times10^{-6}$ & \textbf{99.9\%} \\
\midrule
Totale $V_{\text{noise}}$ & $8.82\times10^{-6}$ & 100\% \\
\midrule
NEP & $2.26\times10^{-6}$ & --- \\
\bottomrule
\end{tabular}
\end{center}

\subsection{Power sweep (120 punti, $N_{\text{MC}} = 1000$)}

\begin{center}
\begin{tabular}{lrl}
\toprule
Metrica & Valore & Note \\
\midrule
NEP & $2.26\times10^{-6}$ W & soglia inferiore del range utile \\
$P_{\text{sat}}$ & $0.256$ W & soglia superiore \\
Punti noise-dom. & 45/120 & $P < \text{NEP}$ \\
Punti ricostruibili & 67/120 & range utile \\
Punti saturi & 8/120 & $P > P_{\text{sat}}$ \\
RMSE (punti ok) & $7.4\times10^{-8}$ W & \\
$\sigma_P$ (MC) vs NEP & $2.30\times10^{-6}$ vs $2.26\times10^{-6}$ W & consistente \\
\bottomrule
\end{tabular}
\end{center}

% ============================================================
\section{Guida all'Estensione}

\subsection{Modificare i parametri fisici}

La strada più semplice è usare l'override in \texttt{bolometer\_params}:
\begin{lstlisting}
p = bolometer_params('R0', 200, 'alpha', 3.4e-3, 'G', 5e-5);
sweep = run_power_sweep(p);
\end{lstlisting}

\subsection{Aggiungere una nuova sorgente di rumore}

Aggiungere il campo in \texttt{bolometer\_params.m}, poi aggiornare
\emph{sia} \texttt{add\_noise\_to\_signal.m} \emph{sia}
\texttt{reconstruct\_power.m} (che ricalcola internamente $V_{\text{noise}}$
per la classificazione dei regimi):

\begin{lstlisting}[caption={Schema per aggiungere una nuova sorgente di rumore}]
% in bolometer_params.m
params.i_noise = 1e-12;    % nuovo: corrente di rumore dell'amplif. [A/rtHz]

% in add_noise_to_signal.m
V_I = params.i_noise * params.R0 * sqrt(params.BW);  % nuova sorgente
V_noise = sqrt(V_J^2 + V_A^2 + V_Q^2 + V_I^2);      % aggiungi qui
\end{lstlisting}

\subsection{Aggiungere dinamica temporale}

Attualmente il modello è in regime stazionario. Per simulare la
risposta temporale, sostituire $\Delta T = P/G$ con l'integrazione ODE:

\begin{lstlisting}[caption={Schema per risposta temporale}]
% in forward_model.m -- modalita' dinamica
% C * d(DeltaT)/dt = P - G * DeltaT
% soluzione analitica per P costante a tratti:
tau = params.C / params.G;
DeltaT_t = (P/params.G) .* (1 - exp(-t/tau));
\end{lstlisting}

\subsection{Aggiungere più configurazioni a confronto}

\begin{lstlisting}[caption={Confronto tra configurazioni in main.m}]
configs = {
    bolometer_params(),                         % default
    bolometer_params('R0', 200),                % R0 doppio
    bolometer_params('G', 5e-5),                % G dimezzata
    bolometer_params('N_bits', 24),             % ADC 24-bit
};
labels = {'default', 'R0=200', 'G=5e-5', 'ADC 24-bit'};

for k = 1:numel(configs)
    opt.do_plot  = false;
    sw{k} = run_power_sweep(configs{k}, opt);
    fprintf('%s:  NEP=%.2e W,  DR=%.1f dB\n', ...
        labels{k}, sw{k}.rec.NEP, ...
        20*log10(sw{k}.fwd.P_sat / sw{k}.rec.NEP));
end
\end{lstlisting}

\subsection{Estendere la sensitivity analysis}

Per aggiungere un nuovo parametro al sweep (es.\ \texttt{N\_bits}):

\begin{lstlisting}[caption={Aggiungere N\_bits al sweep in sensitivity\_analysis.m}]
% nel vettore sweep_defs:
struct('name','N_bits', 'label','N_{bits}', ...
       'base', base_params.N_bits, ...
       'lo_fac', 0.5, 'hi_fac', 2.0)   % da 8 a 32 bit
\end{lstlisting}

\begin{warnbox}
\textbf{Attenzione:} \texttt{N\_bits} è intero --- il range \texttt{logspace}
produrrà valori non interi. Aggiungere \texttt{round()} nella lettura
oppure usare un range lineare con \texttt{linspace}.
\end{warnbox}

% ============================================================
\section{Aggiornamento del Contesto}

Il file \texttt{CONTEXT.md} nella cartella \texttt{bolometer/} deve essere
aggiornato ad ogni modifica rilevante. Serve come memoria rapida per
sessioni di lavoro future senza dover rileggere tutto il codice.

Il formato è:
\begin{itemize}[noitemsep]
  \item Tabella file con stato (\texttt{✅}/\texttt{⬜}) aggiornata
  \item Sezione per ogni blocco con i risultati numerici attesi
  \item Sezione \emph{Come eseguire} con esempi aggiornati
  \item Timestamp in fondo
\end{itemize}

% ============================================================
\section*{Appendice: Parametri e Unità di Misura}

\begin{center}
\begin{tabular}{lll}
\toprule
Simbolo & Definizione & Unità SI \\
\midrule
$P$                  & potenza assorbita             & W \\
$\Delta T$           & incremento di temperatura     & K \\
$\Delta R$           & variazione di resistenza      & $\Omega$ \\
$\Delta V$           & variazione di tensione        & V \\
$R_0$                & resistenza a $T_0$            & $\Omega$ \\
$\alpha$             & TCR (Temperature Coeff.\ Resistance) & K$^{-1}$ \\
$G$                  & conduttanza termica           & W\,K$^{-1}$ \\
$C$                  & capacità termica              & J\,K$^{-1}$ \\
$\tau = C/G$         & costante di tempo termica     & s \\
$I_{\text{bias}}$    & corrente di polarizzazione    & A \\
$S$                  & sensibilità $= I_{\text{bias}}R_0\alpha/G$ & V\,W$^{-1}$ \\
$k_B$                & costante di Boltzmann         & J\,K$^{-1}$ \\
$\Delta f$           & larghezza di banda            & Hz \\
$e_n$                & densità rumore amplif.        & V\,$\sqrt{\text{Hz}}^{-1}$ \\
$N$                  & risoluzione ADC               & bit \\
$V_{\text{range}}$   & full-scale ADC                & V \\
$V_{\text{noise}}$   & rumore totale rms             & V \\
NEP                  & Noise Equivalent Power        & W \\
DR                   & Dynamic Range $= P_{\text{sat}}/\text{NEP}$ & --- (dB) \\
\bottomrule
\end{tabular}
\end{center}

\end{document}
